\subsection*{Slide Examples}
%------------------------------
\begin{myslide}{c}{Bulleted List Slide Example}

\begin{columns}

\begin{column}{1\textwidth}

\begin{itemize}

\item \lipsum[1][1-2]

\item \lipsum[1][3]

\item \lipsum[1][4-5]

\item \lipsum[1][6]

\end{itemize}

\end{column}

\end{columns}
\annotation{
This slide demonstrates a simple bulleted list layout using the myslide environment. The content uses Lorem Ipsum text to show how itemize environments work within the column structure. This is the most basic slide format in the template.
}
\end{myslide}

%------------------------------
\begin{myslide}{c}{Side-by-Side Bulleted List + Figure Slide Example}

\begin{columns}

\begin{column}{0.65\textwidth}

\begin{itemize}

\item \lipsum[1][1-2]

\item \lipsum[1][3]

\item \lipsum[1][4-5]

\item \lipsum[1][6]

\end{itemize}

\end{column}

\myfigcol{\mygraphicspath/csu.pdf}{0.25}

\end{columns}
\annotation{
This slide shows a two-column layout with text on the left and a figure on the right. The myfigcol command is used to place the figure in the right column. This is a common layout for presentations where you want to combine explanatory text with visual content.
}
\end{myslide}

%------------------------------
\begin{myslide}{c}{Top-Bottom Figure + Bulleted List Slide Example}

\myfig{\mygraphicspath/csu.pdf}{0.25}

\myspace

\begin{columns}

\begin{column}{1\textwidth}

\begin{itemize}

\item \lipsum[1][1-2]

\item \lipsum[1][3]

\item \lipsum[1][4-5]

\item \lipsum[1][6]

\end{itemize}

\end{column}

\end{columns}
\annotation{
This slide demonstrates a vertical layout with a figure at the top and bulleted list below. The myfig command centers the figure, and myspace adds appropriate spacing between the figure and the text content. This layout is useful when the figure needs to be prominent and the text provides detailed explanation.
}
\end{myslide}

%------------------------------
\begin{myslide}{c}{Side-by-Side Figures Slide Example}

\begin{columns}

\myfigcol{\mygraphicspath/csu.pdf}{0.25}

\myfigcol{\mygraphicspath/csu.pdf}{0.25}

\end{columns}
\annotation{
This slide shows how to place multiple figures side-by-side using the columns environment and myfigcol commands. Each figure is placed in its own column, allowing for easy comparison or showing related visual content together.
}
\end{myslide}

\subsection*{Other}
%------------------------------
\begin{myslide}{c}{Colors}

\begin{itemize}

\item General colors:
\textcolor{niceblue}{niceblue}
\textcolor{nicered}{nicered}
\textcolor{nicegreen}{nicegreen}
\textcolor{nicepink}{nicepink}
\textcolor{nicepurple}{nicepurple}
\textcolor{nicegray}{nicegray}

\item Grays:
\textcolor{lightgray}{lightgray}
\textcolor{mediumgray}{mediumgray}
\textcolor{darkergray}{darkergray}

\item CSU colors:
\textcolor{csugreen}{csugreen}
\textcolor{csugold}{csugold}
\textcolor{csugreendarker}{csugreendarker}


\item Beamer theme colors (global):
\textcolor{fixedtextcolor}{fixedtextcolor}
\textcolor{fixedbgcolor}{fixedbgcolor}
\textcolor{tcolorboxbg}{tcolorboxbg}
\textcolor{primarycolor}{primarycolor}
\textcolor{graycolor}{graycolor}

\item Beamer theme colors (light/dark):
\textcolor{textcolor}{textcolor}
\textcolor{textcolorlight}{textcolorlight}
\textcolor{textcoloremph}{textcoloremph}
\textcolor{bgcolor}{bgcolor}
\textcolor{bgcolorlight}{bgcolorlight}
\textcolor{edgecolor}{edgecolor}

\item Change the colors in \lstinline{.config/2-colors.tex}

\end{itemize}
\annotation{
This slide displays all available colors in the template. The template includes general colors, grayscale options, CSU-specific colors, and Beamer theme colors that adapt to light and dark modes. All colors can be customized in the 2-colors.tex configuration file.
}
\end{myslide}

%------------------------------
\begin{myslide}{c}{Nested \lstinline{itemize} and \lstinline{enumerate} Environments}

\begin{itemize}

\item Apples

\begin{itemize}

\item Apples

\begin{itemize}

\item Apples

\item Apples

\end{itemize}

\item Apples

\item Apples

\end{itemize}

\item Apples

\end{itemize}

\begin{enumerate}

\item Oranges

\begin{enumerate}

\item Oranges

\item Oranges

\begin{enumerate}

\item Oranges

\item Oranges

\end{enumerate}

\item Oranges

\end{enumerate}

\item Oranges

\item Oranges

\end{enumerate}
\annotation{
This slide demonstrates nested list environments. Both itemize (bullet points) and enumerate (numbered lists) can be nested multiple levels deep. The template automatically handles the indentation and styling for nested lists, making it easy to create hierarchical content structures.
}
\end{myslide}